\documentclass[11pt]{article}

% Change "review" to "final" for the camera-ready version.
\usepackage[review]{acl}

\usepackage{times}
\usepackage{latexsym}
\usepackage[T1]{fontenc}
\usepackage[utf8]{inputenc}
\usepackage{microtype}
\usepackage{inconsolata}
\usepackage{graphicx}
\usepackage{booktabs}
\usepackage{tabularx}
\usepackage{array}
\usepackage{amsmath}
\usepackage{amssymb}

\title{Semantic Similarity in Latent Space for LLM Prediction Evaluation}

\author{
  Team Members: Alice, Bob, Charlie, David \\
  AIAA 4051 Final Research Project \\
  Introduction to Natural Language Processing, HKUST(GZ) \\
}

\begin{document}
\maketitle

\begin{abstract}
This project studies whether embedding-space semantic similarity can be used as an automatic proxy for the correctness of large language model answers in question answering. We generate answers with Qwen2.5-7B-Instruct, encode prediction-reference pairs with MiniLM and BGE embeddings, and evaluate whether cosine similarity separates automatically labeled correct and incorrect predictions. Across four datasets, similarity is highly predictive for short-answer QA, reaching ROC--AUC values of 0.962--0.990 on SciQ and SimpleQuestions-Wiki. However, the original Natural Questions setting fails because short model predictions are compared against long evidence passages, producing inverted or weak similarity signals. Manual failure analysis shows that many errors arise from topical relatedness, answer granularity mismatch, and automatic-label artifacts. We implement an NQ reference-extraction improvement that converts long evidence passages into concise answer references, improving NQ ROC--AUC from 0.269/0.391 to 0.705/0.711 for MiniLM/BGE on the same 500-example subset.
\end{abstract}

\section{Introduction}

Automatic evaluation of free-form QA predictions is difficult because surface-form matching often underestimates correct paraphrases, aliases, and semantically equivalent answers. Embedding similarity offers an appealing alternative: if a predicted answer and a reference answer are close in latent semantic space, the prediction may be treated as correct.

This report evaluates that idea for both short-form and long-form QA. Our central question is:
\begin{quote}
Can cosine similarity between prediction and reference embeddings serve as a reliable proxy for QA correctness?
\end{quote}

The project covers the full workflow: LLM prediction generation, automatic correctness labeling, embedding representation, similarity scoring, empirical evaluation, failure analysis, and an implemented improvement for long-reference datasets.

\section{Task and Pipeline}

Given a question $q$, a reference answer $r$, and an LLM prediction $p$, we compute an embedding similarity score
\begin{equation}
s_m(p,r) =
\frac{E_m(p) \cdot E_m(r)}
{\lVert E_m(p) \rVert \lVert E_m(r) \rVert},
\end{equation}
where $E_m$ is an embedding model. We use two sentence embedding models: MiniLM and BGE, following the broader sentence-embedding evaluation setting introduced by Sentence-BERT-style encoders \citep{reimers-gurevych-2019-sentencebert}.

The pipeline has six stages:
\begin{enumerate}
    \item load QA datasets and references;
    \item generate predictions using Qwen2.5-7B-Instruct;
    \item create rule-based automatic correctness labels;
    \item compute embedding similarities for prediction-reference pairs;
    \item evaluate threshold-based classifiers using F1 and ROC--AUC;
    \item inspect failures and propose improvements.
\end{enumerate}

\subsection{Automatic Correctness Labels}

The field \texttt{correct\_label} is a rule-based automatic label, not a human annotation. It combines exact match, token overlap/F1, and containment-style checks as implemented in \texttt{src/correctness\_labeling.py}. These labels provide a reproducible reference signal for threshold evaluation, but they are imperfect: they can reject valid paraphrases and aliases, especially in open-ended or long-form settings.

\subsection{Similarity Decision}

For each embedding model, a similarity-based classifier predicts correctness by thresholding cosine similarity. We report fixed-threshold F1, best-threshold F1, and ROC--AUC. We also compute the similarity gap:
\begin{equation}
\Delta = \mathbb{E}[s(p,r) \mid y=1] -
\mathbb{E}[s(p,r) \mid y=0],
\end{equation}
where $y$ is the automatic correctness label. A positive gap means automatically correct answers are more similar to references than incorrect answers.

\section{Datasets}

We evaluate on two short-form QA datasets and two long-form/open-ended datasets: SciQ \citep{welbl-etal-2017-sciq}, SimpleQuestions-Wiki \citep{bordes-etal-2015-simplequestions}, Natural Questions \citep{kwiatkowski-etal-2019-natural}, and TruthfulQA \citep{lin-etal-2022-truthfulqa}. Table~\ref{tab:datasets} summarizes the processed evaluation sets.

\begin{table}[t]
\centering
\small
\begin{tabular}{lrrr}
\toprule
Dataset & $N$ & Correct & Ref. tok. \\
\midrule
SciQ & 500 & 60.6\% & 1.5 \\
SimpleQ-Wiki & 500 & 18.4\% & 1.8 \\
Natural Questions & 5000 & 3.7\% & 100.3 \\
TruthfulQA & 500 & 8.2\% & 9.6 \\
\bottomrule
\end{tabular}
\caption{Dataset statistics. Auto-correct rates come from the project's automatic correctness labels, not from human relabeling.}
\label{tab:datasets}
\end{table}

The most important dataset property is reference granularity. SciQ and SimpleQuestions-Wiki usually provide short answer spans, so the compared texts are similar in length and function. Natural Questions (NQ), in contrast, originally compares short predictions with long Wikipedia evidence passages. This mismatch is the main source of failure in the original NQ results.

\section{Experimental Results}

Table~\ref{tab:main-results} reports the main embedding-similarity results. Short-answer datasets show strong separation between correct and incorrect predictions. On SciQ, MiniLM and BGE achieve ROC--AUC scores of 0.962 and 0.964. On SimpleQuestions-Wiki, they reach 0.987 and 0.990.

\begin{table*}[t]
\centering
\small
\begin{tabular}{llrrrrrr}
\toprule
Dataset & Type & MiniLM AUC & BGE AUC & MiniLM F1 & BGE F1 & MiniLM gap & BGE gap \\
\midrule
NQ original & long & 0.276 & 0.381 & 0.072 & 0.072 & -0.157 & -0.042 \\
SciQ & short & 0.962 & 0.964 & 0.901 & 0.913 & 0.455 & 0.268 \\
SimpleQ-Wiki & short & 0.987 & 0.990 & 0.892 & 0.934 & 0.576 & 0.430 \\
TruthfulQA & long & 0.852 & 0.903 & 0.551 & 0.549 & 0.300 & 0.216 \\
NQ extracted & long & 0.705 & 0.711 & 0.282 & 0.286 & 0.195 & 0.128 \\
\bottomrule
\end{tabular}
\caption{Main results. F1 is the best F1 after threshold search. NQ original uses \texttt{results\_nq\_5000}; NQ extracted uses the implemented reference-extraction run \texttt{results\_nq\_500}.}
\label{tab:main-results}
\end{table*}

The original NQ setting is the main negative case. The gap is negative for both embedding models, which means automatically incorrect predictions are, on average, more similar to the long reference passages than automatically correct predictions. This confirms that raw embedding similarity can reward topical overlap rather than answer correctness when the reference is a long evidence passage.

TruthfulQA is more open-ended than SciQ and SimpleQuestions-Wiki, but its average reference length is much shorter than NQ. Its ROC--AUC values remain useful, especially for BGE at 0.903, although best-threshold F1 is lower because the automatic positive class is small.

\section{Failure Analysis}

We define a similarity failure as a disagreement between the thresholded similarity decision and the automatic correctness label. We inspect two directions:
\begin{itemize}
    \item \textbf{High-similarity wrong}: the prediction is automatically labeled incorrect but receives a high similarity score.
    \item \textbf{Low-similarity correct}: the prediction is automatically labeled correct but receives a low similarity score.
\end{itemize}

Manual annotation separates true limitations of embedding similarity from artifacts caused by strict automatic labels. Table~\ref{tab:manual-categories} summarizes the aggregate manual categories.

\begin{table}[t]
\centering
\small
\begin{tabular}{lr}
\toprule
Manual category & Count \\
\midrule
Semantic similarity limitation & 172 \\
Low-similarity false negative & 78 \\
Automatic label artifact & 66 \\
\bottomrule
\end{tabular}
\caption{Aggregate manual annotation categories over sampled failure cases.}
\label{tab:manual-categories}
\end{table}

\subsection{Observed Failure Modes}

\textbf{Topic relatedness from long passages.}
This is dominant in original NQ. The prediction and reference passage are often about the same entity or event, so the embedding score is high even when the prediction does not answer the question correctly.

\textbf{Related but wrong answers.}
Some predictions are semantically close to the gold answer but differ in relation, entity, or required specificity. Embedding similarity captures neighborhood relatedness but not task-specific correctness.

\textbf{Granularity mismatch.}
Under-specific or over-specific answers can receive high similarity, even when QA correctness requires an exact answer type or level of detail.

\textbf{Label artifacts.}
Automatic rules sometimes mark acceptable paraphrases, aliases, date variants, or numeric equivalents as incorrect. These cases are not pure embedding failures; they reveal label noise in the automatic reference signal.

\textbf{Context dilution.}
Correct predictions can receive low similarity when the answer is surrounded by extra generated context, or when the reference is a long passage that dilutes the answer-bearing span.

\section{Implemented Improvement: NQ Reference Extraction}

The current branch implements an NQ-specific reference extraction module in \texttt{src/reference\_answer.py}. Instead of embedding the full evidence passage, the pipeline extracts a shorter answer-like reference before labeling and similarity computation.

Table~\ref{tab:nq-improvement} compares original passage references with extracted references on the same 500-example subset. The improvement is large for both embedding models: MiniLM improves from 0.269 to 0.705 ROC--AUC, and BGE improves from 0.391 to 0.711.

\begin{table}[t]
\centering
\small
\begin{tabular}{lrrrr}
\toprule
Model & AUC$_b$ & AUC$_a$ & Gap$_b$ & Gap$_a$ \\
\midrule
MiniLM & 0.269 & 0.705 & -0.149 & 0.195 \\
BGE & 0.391 & 0.711 & -0.031 & 0.128 \\
\bottomrule
\end{tabular}
\caption{Effect of NQ reference extraction on the same 500-example subset. Subscripts $b$ and $a$ indicate before and after extraction.}
\label{tab:nq-improvement}
\end{table}

This result supports the main diagnosis: similarity evaluation works better when prediction and reference are comparable answer units. Extraction does not solve all remaining errors. After extraction, high-similarity wrong cases still include related-but-wrong answers and cases requiring entailment or answer-type checks.

\section{Improvement Roadmap}

Beyond the implemented NQ extraction module, the failure analysis suggests several future improvements.

\textbf{Human-annotated calibration.}
The manually labeled failure sample can be used as a small calibration set for threshold tuning and label-noise analysis. This directly addresses cases where automatic labels are stricter than human judgment.

\textbf{Answer normalization.}
Preprocessing can canonicalize aliases, dates, numbers, units, casing, punctuation, and answer spans before similarity scoring.

\textbf{Bidirectional extraction.}
NQ extraction currently repairs references. A more general system could extract concise answer candidates from both predictions and references.

\textbf{Hybrid scoring.}
Embedding similarity can be combined with token F1, containment, entity overlap, answer-type agreement, and other structured overlap signals.

\textbf{Dataset-specific thresholds.}
Short-form and long-form datasets have different similarity distributions, so a single global threshold is unlikely to be optimal.

\textbf{NLI or LLM verifier.}
An entailment layer can distinguish semantic relatedness from actual answer correctness, especially for high-similarity wrong cases.

\section{Conclusion}

Embedding similarity is a strong correctness signal when the prediction and reference are comparable answer units. It works especially well for short-form QA, where references are compact answer spans. Its main weakness is not only the embedding model, but the mismatch between semantic relatedness and task-specific correctness. The original NQ setup demonstrates this failure clearly: long evidence passages make incorrect answers appear semantically similar. The implemented NQ reference-extraction module substantially improves the signal, showing that repairing reference granularity is an effective path for long-form QA evaluation.

\section*{Limitations}

The evaluation uses automatic correctness labels, which are reproducible but noisy. The NQ improvement is evaluated on a 500-example subset, while the original NQ baseline uses 5000 examples. Manual annotation is sampled rather than exhaustive. Finally, threshold tuning reports best-threshold F1 on the available evaluation data; a deployed system would need a separate validation set.

\section*{Acknowledgments}

This report is prepared for the AIAA 4051 final research project. The implementation and analysis use the local project pipeline in \texttt{src/}, \texttt{scripts/}, and \texttt{failures\_analysis\_and\_improvement/}.

\bibliography{custom}

\appendix

\section{Reproducibility Notes}

The main scripts are:
\begin{itemize}
    \item \texttt{scripts/01\_generate\_predictions.py};
    \item \texttt{scripts/02\_compute\_similarity.py};
    \item \texttt{scripts/03\_evaluate.py};
    \item \texttt{scripts/04\_visualize\_results.py};
    \item \texttt{scripts/analyze\_part4\_strict.py}.
\end{itemize}

The summary tables used in this report are stored in the \texttt{summary\_tables} subfolder under \texttt{failures\_analysis\_and\_improvement}. The implemented NQ extraction code is in \texttt{src/reference\_answer.py}.

\end{document}
